
%!TeX root = main.tex
%!encoding = UTF-8 Unicode

\chapter{Abstract}

This thesis examines the interface between two domains of power electronics that are frequently treated as distinct despite their fundamental interdependence.
Although semiconductor device switching is the primary source of conducted and radiated emissions, electromagnetic compatibility (EMC) remains predominantly addressed as a system-level problem rather than as an intrinsic semiconductor design variable.
The disconnect stems from a fundamental measurement limitation: conventional characterization methods for power semiconductor devices cannot resolve the high-frequency, EMC-relevant signal components emitted by the device because they fall below the measurement noise floor.
Since a quantity cannot be optimized without first being measurable, the lack of characterization has prevented electromagnetic noise from being treated as an intrinsic semiconductor design target.
Consequently, noise optimization is often deferred to later design stages, where mitigation is implemented through converter- or system-level measures at the expense of efficiency, compactness, and development time.

The thesis contributes a measurement methodology that enables direct, low-noise characterization of intrinsic electromagnetic emissions at the semiconductor power terminals.
While conventional methods cannot resolve the EMC-relevant signals, which are masked by thermal and quantization noise and therefore remain unobservable, the proposed methodology enables their direct measurement and analysis.
The approach is based on a dual-channel concept for both a differential active voltage probe and a current probe, in which a dedicated high-frequency path, optimized independently of the conventional low-frequency measurement path, amplifies the EMC-relevant signal components above the system noise floor.
Residual noise is further suppressed by a digital post-processing chain combining oversampling and averaging with a time-dependent noise filter that concentrates the filtering outside the switching transient, leaving the transient itself unaltered.
Together, these methods enable semiconductor-focused EMC analysis up to \SI{400}{MHz} while preserving conventional performance indicators such as switching losses, overvoltage, and \sddt{v}.
An additional event-resolved spectral analysis method provides deeper insight into individual noise contributions.
Because the methodology renders a previously unobservable behavior measurable for the first time, particular emphasis is placed on verification of the measurement methodology and associated analysis techniques.

Applying the methodology reveals differences in the frequency range of \SIrange{10}{400}{MHz} between SiC MOSFET variants and operating points, quantifying effects that were previously only assumed to be relevant.
It further links switching behavior directly to EMC-relevant spectral components, establishing device emissions as a measurable design target.
The developed approach is not limited to a specific semiconductor device technology and provides a general basis for semiconductor-focused EMC analysis at multiple stages of power electronic system development:
(1) during device technology development, it enables characterization of different design variants and direct evaluation of trade-offs with conventional performance metrics;
(2) during converter design, it enables benchmarking of candidate semiconductor devices with respect to their electromagnetic emissions for a given application; and
(3) during compliance testing, it enables identification of the switching event most likely to dominate emissions and cause violations of the applicable standard limits.
Collectively, these contributions support earlier, EMC-aware design decisions across future power semiconductor and converter development.